% Restructured book: Chapter 1 imported from Chap 1_ Motivation and Objectives.docx.
\chapter{Introduction}

\section{Motivation}

Interior layout design is a practical need in Vietnam's urban housing context. In 2024, the Ministry of Construction reported about 125,545 successful apartment and individual-house transactions \cite{Vietnam2025InformationHousingReal}, showing a large housing market where many households may face furnishing, renovation, or reorganization decisions. After moving into, renting, or reorganizing a room, users often need to decide where beds, wardrobes, sofas, desks, storage units, and media furniture should be placed. These decisions affect cost, circulation, and daily usability. If the initial layout ignores room boundaries, openings, or furniture dimensions, users may need avoidable revisions before purchasing furniture or implementing the arrangement.

Hiring an architect or an interior designer is appropriate when a project requires final drawings, construction documents, detailed material selection, and professional visualization. However, this service usually requires additional cost and time. Interior design services in Vietnam typically cost around 150,000--250,000 VND/m\textsuperscript{2}, with the design process often taking 10--20 days depending on project complexity and revisions \cite{NoiThatMansion2026ApartmentDesignPrice,Hometalk2026ApartmentInteriorDesignPrice}. At the early planning stage, many users simply want to explore different furniture layouts, verify space feasibility, and get a sense of the overall design direction. A lightweight tool that accepts ordinary user descriptions and quickly produces plausible, editable layouts can help users clarify their needs before professional refinement becomes necessary.

This situation leads to the first problem addressed in this thesis: supporting non-professional users in creating an initial furniture layout by themselves. The target users include people arranging their own bedrooms or living rooms, apartment occupants comparing furniture options, and small accommodation owners who need a quick spatial concept before professional refinement. For these users, a useful layout must do more than place objects inside a room. It must respect the room boundary, furniture dimensions, doors, windows, collision avoidance, and basic circulation. It should also remain editable because early requirements are often incomplete and may change after the user sees the first design.

The second problem is visualization. A top-down layout is useful for checking positions and distances, but it does not fully show how the room may feel in practice. From a 2D plan alone, users may still find it difficult to imagine materials, color harmony, lighting, depth, and overall atmosphere. A rendered perspective image helps them evaluate whether the selected layout is visually acceptable before they continue revising the design or purchasing furniture.

These two problems are closely connected. Layout generation provides spatial control and editability, while image rendering provides visual feedback. If the rendered image is generated without respecting the layout, it may look attractive but become difficult to use as a design reference. If the system only provides an editable layout without realistic visualization, the user may still struggle to assess the design. Therefore, this thesis focuses on two connected problems in early-stage interior design: generating an editable and geometrically valid initial furniture layout for non-professional users, and producing realistic visualization that remains consistent with the selected layout.

\section{Objectives and Scope of the Thesis}

The two problems identified in Section 1.1 require more than visual inspiration. Existing AI-based interior design tools can generate attractive room images or suggest plausible object arrangements, but they often do not preserve the full chain from user intent and room geometry to editable layout data and layout-consistent visualization. Image-based tools usually return static pixels, while layout-planning approaches still need geometric validation before their outputs can be used reliably. Even when visual generation is conditioned on a layout, the result may still drift from the selected room structure or furniture arrangement. These limitations motivate the objectives of this thesis.

From these limitations, this thesis defines two main objectives. The first objective is to build a system that can generate an initial furniture layout from an ordinary-language user brief, explicit room geometry, fixed openings, and a furniture inventory. The generated layout must be stored as structured room-and-object data so that it can be checked for basic spatial constraints, displayed in 2D/3D, and manually refined by the user through direct editing. The second objective is to generate realistic perspective images from the selected layout. These images should help users understand the visual direction of the room while preserving the main room structure, camera view, openings, and furniture arrangement defined by the editable layout.

The scope of this thesis is limited to living rooms and bedrooms. These room types are selected because they are common in apartments and contain representative layout constraints. A living room usually involves seating, tables, media furniture, storage, and circulation relationships. A bedroom usually involves bed placement, wardrobe access, bedside furniture, lighting, and optional work or entertainment areas. The system takes as input a natural-language design brief, a room boundary, fixed openings, and a furniture inventory. In this thesis, the inventory is a database of furniture records with category, dimensions, style and material metadata, and optional 3D asset files. The expected outputs are an editable 2D layout, a 3D preview, a generated render image, and controlled image-editing results.

The thesis does not aim to replace architects or professional interior designers. It also does not address construction-level tasks such as structural design, electrical and plumbing systems, technical lighting design, ventilation, cost estimation, or legal construction drawings. The proposed system is intended as an early-stage design assistant. Its purpose is to help users explore possible layouts, revise them interactively, and obtain visual feedback before moving to more detailed design decisions.

\section{Proposed Approach}

To achieve these objectives, the thesis proposes RoomFlow, an AI-assisted room design system that combines language understanding, geometric validation, editable layout data, and layout-based visualization. The key idea is that RoomFlow should not generate a final image directly from a text prompt. Instead, it first builds a structured room layout that can be checked, edited, and previewed. This layout then becomes the basis for 3D preview, perspective rendering, and later image editing. In this way, the proposed solution keeps spatial control in the editable layout stage while using image generation mainly to improve visual realism and user understanding.

For layout generation, the system takes the room type, room boundary, fixed openings, furniture inventory, and user brief as input. A multi-agent pipeline is used to interpret the user’s intent, identify required furniture, and infer high-level spatial preferences. The final placement is then checked by a geometric solver rather than relying on the language model alone. This separation allows the system to validate containment, collisions, opening clearance, and circulation before returning an editable layout that users can review and modify.

For visualization, the selected editable layout is reconstructed as a 3D scene using Three.js. The user can inspect the room, adjust the camera view, and use the selected view as a layout reference for image generation. The render is therefore not generated from text alone, but guided by the room shell, openings, visible furniture, camera view, and preservation constraints. This helps the generated image improve realism while still remaining consistent with the selected layout.

After the initial render is generated, the system also supports controlled image editing. Users can request changes such as adjusting materials, colors, lighting, or small decorative elements. These edits are treated as visual refinements rather than layout changes, so the system should preserve the room structure, camera view, openings, and major furniture positions unless the user explicitly asks otherwise. This keeps image editing useful for exploration without breaking the selected design.

Overall, the proposed solution treats layout data as the central representation of the design. The language model helps convert user intent into design requirements, the validation module ensures that the layout remains spatially usable, the 3D preview helps users inspect the arrangement, and the image generation module provides realistic visual feedback. By connecting these components, the system supports the two main problems of this thesis: creating editable initial layouts and producing layout-consistent visualization for early-stage interior design.


\section{Thesis Organization}

The remainder of this thesis is organized into seven chapters.

Chapter 2 surveys the current status of AI-assisted interior design and analyzes the system requirements. The chapter reviews representative directions, including prompt-based image generation, LLM-based layout planning, hybrid layout generation, and layout-conditioned image generation. From their limitations, the chapter identifies the main requirements of the proposed system and specifies the key use cases, functional requirements, and non-functional requirements.

Chapter 3 establishes the theoretical background and technology foundation of the thesis. The chapter introduces indoor layout representation, geometric constraints for furniture placement, LLM-assisted planning, structured agents, indoor scene synthesis, layout-conditioned image generation, evaluation metrics, and the web technologies used in the system.

Chapter 4 presents the solution for the first problem: automatic editable furniture layout design. The chapter describes the overall processing flow, input representation, semantic planning module, room modeling module, object program construction, geometric solver, and the preparation of editable layout outputs.

Chapter 5 presents the solution for the second problem: snapshot-guided image generation and editing. The chapter explains how the selected layout is converted into visual guidance for perspective rendering, how scene and visible-object constraints are used to preserve layout consistency, and how localized image editing is controlled after the initial render is generated.

Chapter 6 describes the construction of the proposed system. The chapter presents the software architecture, package design, user interface design, backend pipeline, 3D snapshot and rendering design, database design, main implementation tools, implemented functions, and deployment configuration.

Chapter 7 reports the experimental evaluation. The chapter evaluates the proposed system on both layout generation and image generation/editing, using benchmark scenarios, external baselines, ablation settings, quantitative metrics, and representative qualitative examples. It also discusses observed strengths, limitations, and threats to validity.

Chapter 8 concludes the thesis and outlines future work. The chapter summarizes the main results, discusses the remaining limitations, and suggests possible directions for improving layout generation, visualization consistency, image editing, and user interaction.
